\chapter{Decay-Heat Accidents: Three Mile Island and Fukushima}
\label{ch:decayheat-acc}

Chernobyl was a reactivity-stability accident: the fission power ran away. The two
other defining accidents of the nuclear age---Three Mile Island and
Fukushima---were the opposite kind: the reactors shut down correctly, but the
\emph{decay heat} (Chapter~\ref{ch:decay}) could not be removed. Together with
Chernobyl they complete the picture of reactor thermal safety, showing that
stopping fission is necessary but not sufficient.

\section{Three Mile Island (1979)}

At Three Mile Island Unit~2, a minor fault (a stuck-open pressuriser relief valve)
combined with misleading instrumentation led operators to throttle back the
emergency core cooling under the mistaken belief that the system was full of water.
The reactor had scrammed promptly and correctly---there was no reactivity
problem---but decay heat, initially several per cent of the $2700\unit{MW}$ thermal
rating, continued unabated. With cooling reduced, the core boiled dry in part,
overheated, and suffered substantial melting over the following hours. The
containment building, however, held, and off-site releases were minimal.

TMI's lessons were about \emph{human factors and decay heat}: control-room
instrumentation that misled rather than informed; operator training that had not
prepared the crew for the actual scenario; and the inexorability of decay heat. It
transformed operator training, control-room design, and emergency procedures
worldwide---and, seven years before Chernobyl, established that a reactor is a
thermal hazard for days after shutdown.

\section{Fukushima Daiichi (2011)}

The Fukushima Daiichi accident is the starkest demonstration of the decay-heat
imperative. A magnitude-9 earthquake caused the operating reactors to scram
correctly; fission stopped as designed. But the ensuing tsunami flooded the
emergency diesel generators and switchgear, causing a prolonged \emph{station
blackout}---a complete loss of AC power. With no power to drive cooling, decay heat
(Chapter~\ref{ch:decay}) accumulated in the cores over hours and days. The water
boiled away, the fuel overheated, the zirconium cladding reacted with steam to
generate hydrogen, and the hydrogen exploded, breaching containment in three units
and releasing radioactivity.

Fukushima was not a stability accident and not an operator-violation accident; it
was a failure to remove decay heat after a beyond-design-basis natural event
disabled the active cooling systems. Its central lesson reinforced the move toward
\emph{passive} decay-heat removal (Chapter~\ref{ch:inherent}): cooling that needs
no electrical power and no operator action for a long grace period, precisely
because the active systems can be lost.

\section{Three accidents, three lessons}

The three defining accidents form a complete syllabus of reactor thermal safety:
\begin{itemize}[nosep]
\item \textbf{Chernobyl (1986)}---a \emph{reactivity-stability} failure: positive
feedback and a flawed scram drove a prompt-critical excursion. Lesson: inherent
negative feedback and a sound shutdown system are non-negotiable.
\item \textbf{Three Mile Island (1979)}---a \emph{decay-heat plus human-factors}
failure: correct scram, but cooling was throttled and the core melted. Lesson:
instrumentation, training, and respect for decay heat.
\item \textbf{Fukushima (2011)}---a \emph{decay-heat plus station-blackout}
failure: correct scram, but a natural disaster disabled active cooling. Lesson:
passive, power-independent decay-heat removal.
\end{itemize}
Stability (the subject of this book) addresses the first; passive decay-heat
removal addresses the second and third. A reactor that is both inherently stable
\emph{and} passively cooled is safe against all three failure modes---the goal of
modern design.

\keypoint{Not every reactor accident is a stability accident. TMI and Fukushima
shut down correctly but could not remove decay heat---one through human factors,
one through a natural disaster that disabled active cooling. Thermal safety
therefore demands both inherent stability (against excursions) and passive
decay-heat removal (against loss of cooling).}

\section*{Exercises}
\addcontentsline{toc}{section}{Exercises}
\begin{enumerate}[leftmargin=*,itemsep=2pt]
\item Classify TMI, Fukushima, and Chernobyl by failure mode (reactivity vs.\ decay heat).
\item Explain why correct scram did not prevent core damage at TMI and Fukushima.
\item What is a station blackout, and why is it central to Fukushima?
\item Argue why passive decay-heat removal addresses the Fukushima failure mode.
\item State the single lesson each of the three accidents contributes to thermal safety.
\end{enumerate}
